\documentclass{article}

\usepackage[french]{babel}
\usepackage[utf8]{inputenc}
\usepackage{hyperref}
\usepackage{amssymb}
\usepackage{amsmath}
\usepackage{float}
\usepackage{subcaption}

\usepackage{graphicx}
\graphicspath{{img/}}

%%%%%%%%%%%%%%%% Lengths %%%%%%%%%%%%%%%%
\setlength{\textwidth}{15.5cm}
\usepackage[a4paper, total={6in, 8.5in}]{geometry}

%%%%%%%%%%%%%%%% Variables %%%%%%%%%%%%%%%%
\def\projet{0}
\def\titre{Compression d'image à travers la factorisation SVD}
\def\groupe{0}
\def\equipe{0}
\def\responsible{Wissal Faouzi}
\def\secretary{Prenom Nom}
\def\others{Benjamin Chew, Yani haouchine}

\begin{document}

%%%%%%%%%%%%%%%% Header %%%%%%%%%%%%%%%%
\noindent\begin{minipage}{0.98\textwidth}
  \vskip 0mm
  \noindent
  { \begin{tabular}{p{7.5cm}}
      {\bfseries \sffamily
        Projet \projet} \\
      {\itshape \titre}
    \end{tabular}}
  \hfill
  \fbox{\begin{tabular}{l}
      {~\hfill \bfseries \sffamily Groupe \groupe\ - Equipe \equipe
        \hfill~} \\[2mm]
      Responsable : \responsible \\
      Secrétaire : \secretary \\
      Codeurs : \others
    \end{tabular}}
  \vskip 4mm ~

  ~~~\parbox{0.95\textwidth}{\small \textit{Résumé~:} \sffamily Le but de ce projet consiste à programmer un algorithme permettant de faire de la compression d’images en utilisant des techniques matricielles basée sur la factorisation SVD. Ce type d’algorithme est à relier aux algorithmes de compression avec pertes, dont le plus connu est certainement l’algorithme de compression JPEG, lui-même basé usuellement sur la Discrete Cosine Transform (DCT), une transformation voisine de la transformée de Fourier discrète. }
  \vskip 1mm ~
\end{minipage}

%%%%%%%%%%%%%%%% Main part %%%%%%%%%%%%%%%%
\section{Partie 1} \label{part1}

Test

\section{Partie 2} \label{part2}

\section{Partie 3} \label{part3}

\section{Partie 4} \label{part4}

\bibliographystyle{unsrt}
\bibliography{references}

\end{document}
